\section{The VRA-Aware MCMC Algorithm}
\label{sec:algorithm}

\subsection{Motivation: VRA as a Hard Constraint}

VRA compliance is a legal constraint, not a probabilistic preference.
A plan that reduces a protected district's minority VAP below $50\%$ is legally
impermissible in a Gingles-covered state, regardless of how compact or
population-balanced it may be.
Enforcing VRA as a soft penalty in the Metropolis-Hastings ratio would allow
non-compliant plans to appear in the ensemble, violating the conditional
sampling goal.

The correct approach is a hard filter: any proposed step that would produce a
non-compliant plan is rejected unconditionally.
This reduces the effective state space of the chain to
$\Pi_{\mathrm{vra}}(G, k, \Prot) \subseteq \Pi(G, k)$: the set of valid
plans in which every district $d \in \Prot$ has
$\vap(\pi^{-1}(d), \pi) \geq \theta$.
\vr{} is a Markov chain on $\Pi_{\mathrm{vra}}(G, k, \Prot)$, with the
same long-run mixing properties as \fr{} on the restricted space.

\subsection{Formal Algorithm}

\begin{algorithm}
\caption{\vr{} Chain}
\label{alg:vr}
\begin{algorithmic}[1]
\Require Initial plan $\pi_0$, steps $S$, base seed $b$, minority VAP
         fractions $\mathbf{m} \in [0,1]^n$, threshold $\theta$
\Ensure  Final plan $\pi^*$ and three-way step record
\State $\Prot \gets \{d \in [k] : \vap(\pi_0^{-1}(d), \pi_0) \geq \theta\}$
       \Comment{protected set, fixed at construction}
\State $\pi \gets \pi_0$; \quad
       $N_{\mathrm{acc}} \gets 0$; \quad
       $N_{\mathrm{vra}} \gets 0$; \quad
       $N_{\mathrm{mh}} \gets 0$
\For{$s \gets 0$ \textbf{to} $S - 1$}
  \State $\mathrm{rng\_fwd} \gets \mathrm{SmallRng::seed\_from\_u64}(\mathrm{forward\_seed}(b, s, 0))$
  \State $\mathrm{rng\_rev} \gets \mathrm{SmallRng::seed\_from\_u64}(\mathrm{reverse\_seed}(b, s, 0))$
  \State \Comment{--- Run inner ForestRecomChain proposal ---}
  \State $(d_i, d_j) \gets \mathrm{SelectAdjacentPair}(\pi, \mathrm{rng\_fwd})$
  \State $R \gets \pi^{-1}(d_i) \cup \pi^{-1}(d_j)$
  \State $T_{\mathrm{fwd}} \gets \mathrm{WilsonUST}(G[R], \mathrm{rng\_fwd})$
  \State $c_{\mathrm{fwd}} \gets |\{e \in T_{\mathrm{fwd}} : \text{balanced cut}\}|$
  \If{$c_{\mathrm{fwd}} = 0$} $N_{\mathrm{mh}} \mathrel{+}= 1$; \textbf{continue} \EndIf
  \State $e^* \gets \mathrm{Uniform}(C(T_{\mathrm{fwd}}), \mathrm{rng\_fwd})$
  \State $\pi' \gets \pi$ with district pair $(d_i, d_j)$ split by $e^*$
  \State $T_{\mathrm{rev}} \gets \mathrm{WilsonUST}(G[R], \mathrm{rng\_rev})$
  \State $c_{\mathrm{rev}} \gets |\{e \in T_{\mathrm{rev}} : \text{balanced cut under } \pi'\}|$
  \If{$c_{\mathrm{rev}} = 0$} $N_{\mathrm{mh}} \mathrel{+}= 1$; \textbf{continue} \EndIf
  \State $\alpha_{\mathrm{mh}} \gets \min(1,\; c_{\mathrm{fwd}} / c_{\mathrm{rev}})$
  \If{$\mathrm{Uniform}([0,1)) \geq \alpha_{\mathrm{mh}}$}
        \Comment{MH rejection}
    \State $N_{\mathrm{mh}} \mathrel{+}= 1$; \textbf{continue}
  \EndIf
  \State \Comment{--- VRA hard check (fires after MH acceptance) ---}
  \State $\mathrm{vra\_ok} \gets \mathbf{true}$
  \For{$d \in \Prot$}
    \If{$\vap(\pi'^{-1}(d), \pi') < \theta$}
      \State $\mathrm{vra\_ok} \gets \mathbf{false}$; \textbf{break}
    \EndIf
  \EndFor
  \If{$\neg\,\mathrm{vra\_ok}$}
        \Comment{VRA rejection: revert and count}
    \State $N_{\mathrm{vra}} \mathrel{+}= 1$; \textbf{continue}
  \EndIf
  \State $\pi \gets \pi'$; \quad $N_{\mathrm{acc}} \mathrel{+}= 1$
        \Comment{accept}
\EndFor
\State \textbf{assert} $N_{\mathrm{acc}} + N_{\mathrm{vra}} + N_{\mathrm{mh}} = S$
\State \Return $(\pi,\; N_{\mathrm{acc}},\; N_{\mathrm{vra}},\; N_{\mathrm{mh}})$
\end{algorithmic}
\end{algorithm}

\paragraph{VRA check ordering.}
The VRA check fires \emph{after} the MH acceptance decision (line 20--26 of
Algorithm~\ref{alg:vr}).
This ordering is deliberate: the MH ratio is computed without knowledge of
VRA status, so the proposal sequence is identical to \fr{} given the same seed.
A verifier who replays the chain without VRA checking will see the same
proposals in the same order; the VRA rejections are simply additional
unconditional rejections layered on top.

An alternative ordering---checking VRA before the MH step and only computing
the reverse tree if the proposal is VRA-compliant---would save computation
when VRA rejects.
We adopt the post-MH ordering to maintain proposal-sequence auditability and
to match the spec~\citep{dellaLibera2026vraSection}, accepting the minor
inefficiency of computing $T_{\mathrm{rev}}$ for subsequently VRA-rejected
proposals.

\paragraph{Protected set semantics.}
The protected set $\Prot$ is computed from the initial plan $\pi_0$ at chain
construction and is never updated.
This matches the legal standard: the Gingles obligation is tied to the
submitted plan's majority-minority districts, not to the current chain state.
If a district's minority VAP in the initial plan is exactly at threshold,
it is protected and must remain at or above threshold throughout the chain.

\paragraph{Degenerate case: no protected districts.}
If $\Prot = \emptyset$ (no district in $\pi_0$ has minority VAP $\geq \theta$),
the VRA check never fires and $N_{\mathrm{vra}} = 0$ for all steps.
The chain reduces to standard \fr{}; a CLI warning is emitted to alert the
user that VRA enforcement is a no-op.

\subsection{Seed Specification}

\vr{} inherits the \fr{} seed derivation unchanged:
\begin{align*}
  \mathrm{forward\_seed}(b, s, 0) &= \mathrm{SHA256}\!\left(
    \texttt{"FR\_FORWARD\_"} \,\|\, s^{\mathrm{le64}} \,\|\,
    \texttt{"\_"} \,\|\, 0^{\mathrm{le32}} \,\|\,
    \texttt{"\_"} \,\|\, b^{\mathrm{le64}}
  \right)[0{:}8], \\
  \mathrm{reverse\_seed}(b, s, 0) &= \mathrm{SHA256}\!\left(
    \texttt{"FR\_REVERSE\_"} \,\|\, s^{\mathrm{le64}} \,\|\,
    \texttt{"\_"} \,\|\, 0^{\mathrm{le32}} \,\|\,
    \texttt{"\_"} \,\|\, b^{\mathrm{le64}}
  \right)[0{:}8].
\end{align*}
The prefix constants \texttt{FR\_FORWARD\_} and \texttt{FR\_REVERSE\_} are
version-locked by an L0 test that asserts equality with the constants in
\texttt{ForestRecomChain}: any rename in the parent will fail the test,
preventing silent divergence.

This inheritance has an important consequence for auditability: given the same
base seed $b$, \vr{} and \fr{} produce \emph{identical proposals} at every
step.
Every plan accepted by \vr{} is a plan that \fr{} would also have accepted
(conditional on MH acceptance), and is additionally guaranteed to be
VRA-compliant.
A verifier can confirm this by replaying the chain with VRA checking disabled
and observing that every \vr{}-accepted plan appears in the unconstrained
trajectory.

\subsection{Three-Way Step Accounting}

Each step falls into exactly one of three categories:

\begin{description}
  \item[Accepted ($N_{\mathrm{acc}}$).]
        The proposal passes both MH acceptance and the VRA check.
        The chain moves to $\pi'$.
  \item[MH-rejected ($N_{\mathrm{mh}}$).]
        The proposal is rejected by the Metropolis-Hastings criterion
        (including zero-cut failures).
        The VRA check is not reached.
        The chain stays at $\pi$.
  \item[VRA-rejected ($N_{\mathrm{vra}}$).]
        The proposal passes MH acceptance but fails the VRA check.
        The chain stays at $\pi$.
\end{description}

The identity $N_{\mathrm{acc}} + N_{\mathrm{vra}} + N_{\mathrm{mh}} = S$
holds exactly (enforced by assertion in the implementation).
This three-way decomposition enables precise audit: the VRA rejection rate
$\rho_{\mathrm{vra}} = N_{\mathrm{vra}} / S$ quantifies how often the chain
encounters compliant-to-non-compliant transitions, and the effective sample
size is $N_{\mathrm{acc}}$.

\begin{proposition}[VRA invariance]
\label{prop:vra-invariant}
If the initial plan $\pi_0$ satisfies $\vap(\pi_0^{-1}(d), \pi_0) \geq \theta$
for all $d \in \Prot$, then every plan $\pi_s$ visited by \vr{} satisfies
the same condition for all $s \geq 0$.
\end{proposition}

\begin{proof}
By construction.
The chain starts at $\pi_0 \in \Pi_{\mathrm{vra}}$.
At each step, the VRA check accepts the proposal $\pi'$ only if
$\vap(\pi'^{-1}(d), \pi') \geq \theta$ for all $d \in \Prot$.
If the check fails, the chain stays at $\pi$, which by inductive hypothesis
is in $\Pi_{\mathrm{vra}}$.
Hence $\Pi_{\mathrm{vra}}$ is absorbing under \vr{}, and VRA compliance is
a chain invariant.
\end{proof}

\begin{remark}[Connectivity of $\Pi_{\mathrm{vra}}$]
The above distributional analysis assumes that $\Pi_{\mathrm{vra}}$ is connected under \fr{} moves --- i.e., any two VRA-compliant plans can be connected through a sequence of VRA-compliant intermediate steps. For most redistricting states, this holds because majority-minority districts are surrounded by other minority-concentration tracts, allowing boundary adjustments while maintaining VAP fractions above threshold. However, for highly geographically constrained states where the majority-minority district's minority population is concentrated in a small area with no alternative arrangement, $\Pi_{\mathrm{vra}}$ may contain isolated components. In such cases, VRA-Aware MCMC samples one connected component; practitioners should verify the initial plan is not in an isolated component by checking whether any boundary flip preserves VRA compliance.
\end{remark}

\subsection{Rust Implementation}

The \texttt{VraRecomChain} struct in \texttt{redist-ensemble/src/vra\_recom.rs}:

\begin{verbatim}
pub struct VraRecomChain {
    pub inner: ForestRecomChain,
    pub minority_vap: Vec<f64>,          // per-tract minority VAP fraction
    pub vap_threshold: f64,              // default 0.50
    pub protected_districts: HashSet<u32>,
    pub vra_rejections: u64,
    pub steps_accepted: u64,
    pub steps_taken: u64,
}
\end{verbatim}

The \texttt{step()} method returns a \texttt{VraStepRecord}:

\begin{verbatim}
pub struct VraStepRecord {
    pub accepted: bool,
    pub vra_rejected: bool,   // true if rejected by VRA (after MH)
    pub log_w: f64,           // f64::NAN if vra_rejected or mh_rejected
}
\end{verbatim}

The compositor variant is \texttt{SeedCompositor::VraRecom \{ steps, p, vap\_threshold \}}.

\paragraph{CLI invocation.}
\begin{verbatim}
redist state --state NC --year 2020 \
    --search vra-recom --seeds 1000 --percentile 0.0 \
    --vra-threshold 0.50 --weights-override vra-aligned
\end{verbatim}
